\documentclass[../../../uem_tok_q.tex]{subfiles}

\begin{document}
    座標空間内の4点 $\mathrm{O}(0,\:0,\:0)$，$\mathrm{B}(2,\:0,\:0)$，%
    $\mathrm{C}(1,\:1,\:1)$，$\mathrm{D}(1,\:2,\:3)$ を考える．

    \begin{enumerate}
        \item $\overrightarrow{\mathrm{OP}}\perp\overrightarrow{\mathrm{OA}}$，%
            $\overrightarrow{\mathrm{OP}}\perp\overrightarrow{\mathrm{OB}}$，%
            $\overrightarrow{\mathrm{OP}}\cdot\overrightarrow{\mathrm{OC}}=1$ を満たす%
            点Pの座標を求めよ．
        \item 点Pから直線ABに垂線を下ろし，その垂線と直線ABとの交点をHとする．%
            $\overrightarrow{\mathrm{OH}}$ を $\overrightarrow{\mathrm{OA}}$ と %
            $\overrightarrow{\mathrm{OB}}$ を用いて表せ．
        \item 点Qを $\overrightarrow{\mathrm{OQ}}=\myfraca{3}{4}\overrightarrow{\mathrm{OA}}%
            +\overrightarrow{\mathrm{OP}}$ により定め，Qを中心とする半径$r$の球面$S$を考える．%
            $S$が三角形OHBと共有点を持つような$r$の範囲を求めよ．ただし，%
            三角形OHBは3点O，H，Bを含む平面内にあり，周とその内部からなるものとする．
    \end{enumerate}
\end{document}
